% topic_12.tex — Content only. Compiled via main.tex or standalone wrapper.
% ── No \documentclass, \begin{document} or \end{document} here ──

\chapteropener{12}{Motion in a Circle}{%
  \item Angular Displacement and the Radian
  \item Angular Speed and Period
  \item Centripetal Acceleration
  \item Centripetal Force
  \item Examples of Circular Motion
}

% ===== SECTION 1: ANGULAR DISPLACEMENT =====
\section{Angular Displacement and the Radian}

\begin{definitionbox}{The Radian}
The \textbf{radian} (rad) is the SI unit of angle. One radian is the angle subtended at the centre of a circle by an arc whose length equals the radius.
$$\theta\ (\text{rad}) = \frac{\text{arc length}}{\text{radius}} = \frac{s}{r}$$
\end{definitionbox}

\vspace{0.3em}
\begin{center}
\begin{tikzpicture}[scale=1.1]
  % Circle
  \draw[darkblue, thick] (0,0) circle (1.8);
  % Radii
  \draw[darkblue, thick] (0,0) -- (1.8,0);
  \draw[darkblue, thick] (0,0) -- ({1.8*cos(57.3)},{1.8*sin(57.3)});
  % Arc
  \draw[pinkframe, very thick] (1.8,0) arc (0:57.3:1.8);
  % Labels
  \node[darkblue, font=\small] at (1.0,-0.2) {$r$};
  \node[darkblue, font=\small] at ({0.6*cos(28.6)+0.15},{0.6*sin(28.6)}) {$\theta$};
  \node[pinkframe, font=\small] at (2.4,0.9) {$s = r\theta$};
  % Dot at centre
  \fill[darkblue] (0,0) circle (2pt);
  \node[darkgray, font=\scriptsize, anchor=north] at (0,-0.15) {$1\ \text{rad when}\ s = r$};
\end{tikzpicture}
\end{center}

\begin{keybox}{Key Angle Conversions}
\renewcommand{\arraystretch}{1.4}
\begin{tabular}{@{}ll@{}}
Full circle: & $360° = 2\pi\ \text{rad}$ \\
Half circle: & $180° = \pi\ \text{rad}$ \\
Quarter circle: & $90° = \pi/2\ \text{rad}$ \\
To convert degrees to radians: & multiply by $\pi/180$ \\
To convert radians to degrees: & multiply by $180/\pi$ \\
\end{tabular}
\end{keybox}

% ===== SECTION 2: ANGULAR SPEED =====
\section{Angular Speed and Linear Speed}

\begin{definitionbox}{Angular Speed}
The \textbf{angular speed} $\omega$ of an object moving in a circle is the rate of change of angular displacement:
$$\omega = \frac{\Delta\theta}{\Delta t} \qquad \text{units: rad s}^{-1}$$
For uniform circular motion (constant speed), $\omega$ is constant.
\end{definitionbox}

\begin{formulabox}{Angular Speed and Period}
\begin{center}
\Large $\omega = \dfrac{2\pi}{T} = 2\pi f$
\end{center}
\vspace{0.3em}
\begin{tabular}{@{}ll@{}}
$\omega$ & = angular speed (rad s$^{-1}$) \\
$T$ & = period — time for one complete revolution (s) \\
$f$ & = frequency (Hz $= $ s$^{-1}$) \\
\end{tabular}
\end{formulabox}

\begin{formulabox}{Linear Speed and Angular Speed}
\begin{center}
\Large $v = r\omega$
\end{center}
\vspace{0.3em}
\begin{tabular}{@{}ll@{}}
$v$ & = linear (tangential) speed (m s$^{-1}$) \\
$r$ & = radius of the circular path (m) \\
$\omega$ & = angular speed (rad s$^{-1}$) \\
\end{tabular}
\vspace{0.4em}

Although $\omega$ is constant for uniform circular motion, the \textbf{velocity} is not — its direction changes continuously. The speed $v$ is constant but the direction of motion is always \textbf{tangential} to the circle.
\end{formulabox}

\begin{warningbox}{Speed vs Velocity in Circular Motion}
In uniform circular motion:
\begin{itemize}[itemsep=2pt]
  \item \textbf{Speed} is constant — the magnitude of velocity does not change.
  \item \textbf{Velocity} is not constant — its direction changes at every point.
  \item Because velocity changes, the object is \textbf{accelerating}, even though its speed is constant.
\end{itemize}
\end{warningbox}

% ===== SECTION 3: CENTRIPETAL ACCELERATION =====
\section{Centripetal Acceleration}

\begin{definitionbox}{Centripetal Acceleration}
For an object moving in a circle at constant speed, the acceleration is directed \textbf{towards the centre} of the circle. This is called \textbf{centripetal acceleration}.
\begin{itemize}[itemsep=2pt]
  \item It arises because the \textbf{direction} of velocity is continuously changing.
  \item A force of \textbf{constant magnitude} that is always \textbf{perpendicular to the velocity} produces centripetal acceleration.
  \item Because the force is perpendicular to motion, it does \textbf{no work} — kinetic energy (and hence speed) remains constant.
\end{itemize}
\end{definitionbox}

\begin{formulabox}{Centripetal Acceleration}
\begin{center}
\Large $a = r\omega^2 = \dfrac{v^2}{r}$
\end{center}
\vspace{0.3em}
\begin{tabular}{@{}ll@{}}
$a$ & = centripetal acceleration, directed towards centre (m s$^{-2}$) \\
$r$ & = radius of circular path (m) \\
$\omega$ & = angular speed (rad s$^{-1}$) \\
$v$ & = linear speed (m s$^{-1}$) \\
\end{tabular}
\end{formulabox}

\vspace{0.5em}
\begin{center}
\begin{tikzpicture}[scale=1.1]
  % Circle
  \draw[darkblue, dashed, thick] (0,0) circle (2.0);
  \fill[darkblue] (0,0) circle (2pt);

  % Object at top of circle
  \fill[pinkframe] (0,2.0) circle (5pt);

  % Velocity vector (tangential, to the right)
  \draw[->, >=stealth, midblue, very thick] (0,2.0) -- (1.4,2.0)
    node[right, font=\small, midblue] {$v$ (tangential)};

  % Acceleration vector (towards centre, downward)
  \draw[->, >=stealth, pinkframe, very thick] (0,2.0) -- (0,1.0)
    node[right, font=\small, pinkframe] {$a$ (centripetal)};

  % Centre label
  \node[darkblue, font=\small, anchor=north] at (0,-0.15) {centre};

  % Radius
  \draw[darkgray, thin] (0,0) -- (0,2.0);
  \node[darkgray, font=\small, anchor=east] at (-0.1,1.0) {$r$};
\end{tikzpicture}
\end{center}

% ===== SECTION 4: CENTRIPETAL FORCE =====
\section{Centripetal Force}

\begin{definitionbox}{Centripetal Force}
The \textbf{centripetal force} is the resultant force directed towards the centre of the circle that causes centripetal acceleration. It is not a new type of force — it is whatever force (or combination of forces) acts towards the centre in a given situation.
\end{definitionbox}

\begin{formulabox}{Centripetal Force}
Applying Newton's second law ($F = ma$) with $a = r\omega^2 = v^2/r$:
\begin{center}
\Large $F = mr\omega^2 = \dfrac{mv^2}{r}$
\end{center}
\vspace{0.3em}
\begin{tabular}{@{}ll@{}}
$F$ & = centripetal force, directed towards centre (N) \\
$m$ & = mass of the object (kg) \\
$r$ & = radius of circular path (m) \\
$\omega$ & = angular speed (rad s$^{-1}$) \\
$v$ & = linear speed (m s$^{-1}$) \\
\end{tabular}
\end{formulabox}

\begin{keybox}{What Provides the Centripetal Force?}
The centripetal force is provided by different physical forces depending on the situation:
\renewcommand{\arraystretch}{1.4}
\begin{tabular}{@{}p{5.5cm} p{6cm}@{}}
\textbf{Situation} & \textbf{Force providing centripetal force} \\
\hline
Planet/satellite orbiting a star & Gravitational attraction \\
Car turning on a flat road & Friction between tyres and road \\
Ball on a string, horizontal circle & Tension in the string \\
Electron orbiting nucleus (Bohr model) & Electrostatic (Coulomb) attraction \\
Object on the inside of a curved loop & Normal contact force \\
\end{tabular}
\end{keybox}

\begin{warningbox}{``Centrifugal Force'' is Not Real}
There is no outward ``centrifugal force'' acting on the object. In the reference frame of the ground (inertial frame), the only horizontal force on an object in circular motion is the \textbf{inward} centripetal force. The sensation of being ``pushed outward'' is the result of inertia — the body's tendency to continue in a straight line.
\end{warningbox}

% ===== SECTION 5: EXAMPLES OF CIRCULAR MOTION =====
\section{Examples of Circular Motion}

\subsection{Conical Pendulum}

A mass on a string of length $L$ makes angle $\theta$ with the vertical, moving in a horizontal circle of radius $r = L\sin\theta$.

\begin{formulabox}{Conical Pendulum}
Resolving forces:
\begin{align*}
\text{Vertical:} \quad & T\cos\theta = mg \\
\text{Horizontal (centripetal):} \quad & T\sin\theta = \frac{mv^2}{r} = mr\omega^2
\end{align*}
Dividing: $\tan\theta = \dfrac{r\omega^2}{g} = \dfrac{v^2}{rg}$
\end{formulabox}

\subsection{Car on a Banked Track}

On a banked track (angle $\theta$), the horizontal component of the normal force provides the centripetal force, reducing the need for friction.

\begin{formulabox}{Ideal Banking Angle}
For a vehicle travelling at speed $v$ on a track banked at angle $\theta$, with no friction:
$$\tan\theta = \frac{v^2}{rg}$$
\end{formulabox}

\subsection{Vertical Circular Motion}

For an object on the inside of a vertical loop of radius $r$ at the \textbf{top} of the loop:
$$mg + N = \frac{mv^2}{r}$$
The minimum speed for the object to maintain contact: $N \geq 0 \Rightarrow v_{\min} = \sqrt{gr}$.

% ===== SECTION 6: FORMULA SUMMARY =====
\section{Formula Summary Sheet}

\begin{minipage}{\linewidth}
\begin{tcolorbox}[colback=lightgold, colframe=accentgold]
\renewcommand{\arraystretch}{1.8}
\begin{tabular}{@{}lll@{}}
\toprule
\textbf{Formula} & \textbf{Quantity} & \textbf{Units} \\
\midrule
$\theta = s/r$ & Angular displacement & rad \\
$\omega = \Delta\theta/\Delta t$ & Angular speed & rad s$^{-1}$ \\
$\omega = 2\pi/T = 2\pi f$ & Angular speed from period & rad s$^{-1}$ \\
$v = r\omega$ & Linear speed & m s$^{-1}$ \\
$a = r\omega^2 = v^2/r$ & Centripetal acceleration & m s$^{-2}$ \\
$F = mr\omega^2 = mv^2/r$ & Centripetal force & N \\
\bottomrule
\end{tabular}
\end{tcolorbox}

\vspace{0.5em}
\begin{tcolorbox}[colback=lightblue, colframe=darkblue]
\textbf{Useful:}\quad $2\pi\ \text{rad} = 360°$;\quad
$1\ \text{revolution} = 2\pi\ \text{rad}$;\quad
$v = r\omega$ links linear and angular quantities.
\end{tcolorbox}
\end{minipage}

% ===== SECTION 7: EXAM TECHNIQUE =====
\section{Exam Technique and Problem-Solving Strategy}

\begin{keybox}{Step-by-Step Strategy}
\begin{enumerate}[itemsep=3pt]
  \item \textbf{Find $\omega$}: use $\omega = 2\pi/T$ or $\omega = 2\pi f$ — convert rpm or revolutions per second to rad s$^{-1}$ first.
  \item \textbf{Find $v$}: use $v = r\omega$ if needed.
  \item \textbf{Identify the centripetal force}: decide which physical force (tension, gravity, normal force, friction) provides it.
  \item \textbf{Apply $F = mr\omega^2$ or $F = mv^2/r$}: equate to the expression for that force and solve.
\end{enumerate}
\end{keybox}

\begin{warningbox}{Common Errors — Avoid These!}
\begin{itemize}[itemsep=3pt]
  \item Using \textbf{degrees} instead of radians in $v = r\omega$ or $a = r\omega^2$ — $\omega$ must always be in rad s$^{-1}$.
  \item Confusing \textbf{period} $T$ with frequency $f$ — remember $T = 1/f$.
  \item Forgetting to \textbf{identify the centripetal force} physically — in a free-body diagram, only real forces appear; there is no ``centripetal force'' arrow separate from, say, tension or gravity.
  \item In vertical circle problems, failing to account for \textbf{the component of gravity} that contributes to (or subtracts from) the centripetal force.
  \item Confusing \textbf{$r$} (radius) with \textbf{diameter} — always halve the diameter if that is what is given.
\end{itemize}
\end{warningbox}

% ===== SECTION 8: WORKED EXAMPLES =====
\section{Worked Examples}

\subsection*{Example 1 — Angular and Linear Speed}

\textbf{Question:} A CD rotates at $500\ \text{rpm}$. Calculate (a) the angular speed, and (b) the linear speed of a point $6.0\ \text{cm}$ from the centre.

\begin{examplebox}{Solution}
\textbf{Solution:}\\
(a) Convert rpm to rad s$^{-1}$:
$$\omega = \frac{500 \times 2\pi}{60} = \mathbf{52.4\ \text{rad s}^{-1}}$$

(b) $v = r\omega = 0.060 \times 52.4 = \mathbf{3.14\ \text{m s}^{-1}}$
\end{examplebox}

\subsection*{Example 2 — Centripetal Force on a Car}

\textbf{Question:} A car of mass $1200\ \text{kg}$ travels at $18\ \text{m s}^{-1}$ around a flat circular bend of radius $45\ \text{m}$. Calculate the centripetal force required and state what provides it.

\begin{examplebox}{Solution}
\textbf{Solution:}\\
$$F = \frac{mv^2}{r} = \frac{1200 \times 18^2}{45} = \frac{1200 \times 324}{45} = \mathbf{8640\ \text{N}}$$

This force is provided by \textbf{friction} between the tyres and the road surface, acting towards the centre of the bend.
\end{examplebox}

\subsection*{Example 3 — Satellite Orbit}

\textbf{Question:} A satellite orbits Earth at a radius of $7.5\times10^6\ \text{m}$. The gravitational field strength at this altitude is $7.1\ \text{N kg}^{-1}$. Calculate the orbital period.

\begin{examplebox}{Solution}
\textbf{Solution:}\\
Gravitational force provides centripetal force:
$$mg = mr\omega^2 \quad \Rightarrow \quad \omega = \sqrt{\frac{g}{r}} = \sqrt{\frac{7.1}{7.5\times10^6}} = 9.73\times10^{-4}\ \text{rad s}^{-1}$$

Period:
$$T = \frac{2\pi}{\omega} = \frac{2\pi}{9.73\times10^{-4}} = 6454\ \text{s} \approx \mathbf{1.79\ \text{hours}}$$
\end{examplebox}

% ===== SECTION 9: PRACTICE QUESTIONS =====
\section{Practice Exam Questions}

\subsection*{Section A — Short Answer Questions}

\begin{tcolorbox}[colback=white, colframe=midblue, breakable]

\begin{minipage}{\linewidth}
\textbf{Q1.} Define the radian.
\par\hfill\textit{[1 mark]}
\vspace{2.0cm}
\end{minipage}

\begin{minipage}{\linewidth}
\textbf{Q2.} A wheel of radius $0.35\ \text{m}$ completes $120$ revolutions per minute. Calculate (a) the angular speed and (b) the linear speed of a point on the rim.
\par\hfill\textit{[3 marks]}
\vspace{3.5cm}
\end{minipage}

\begin{minipage}{\linewidth}
\textbf{Q3.} Explain why an object moving in a circle at constant speed is accelerating, and state the direction of this acceleration.
\par\hfill\textit{[2 marks]}
\vspace{3.0cm}
\end{minipage}

\begin{minipage}{\linewidth}
\textbf{Q4.} A stone of mass $0.15\ \text{kg}$ is attached to a string and swung in a horizontal circle of radius $0.80\ \text{m}$ at $3.0\ \text{rev s}^{-1}$. Calculate the tension in the string. (Assume the string is horizontal.)
\par\hfill\textit{[3 marks]}
\vspace{3.5cm}
\end{minipage}

\begin{minipage}{\linewidth}
\textbf{Q5.} Explain why the centripetal force does no work on an object moving in a circle at constant speed.
\par\hfill\textit{[2 marks]}
\vspace{2.5cm}
\end{minipage}

\end{tcolorbox}

\subsection*{Section B — Longer Structured Questions}

\begin{tcolorbox}[colback=white, colframe=midblue, breakable]

\begin{minipage}{\linewidth}
\textbf{Q6.} A car of mass $950\ \text{kg}$ travels over the top of a hill that has a circular cross-section of radius $120\ \text{m}$.
\end{minipage}

\begin{enumerate}[label=(\alph*)]
  \item \begin{minipage}[t]{\linewidth}
  Draw a free-body diagram for the car at the top of the hill, showing and labelling the forces acting.
  \par\hfill\textit{[2 marks]}
  \vspace{4.5cm}
  \end{minipage}

  \item \begin{minipage}[t]{\linewidth}
  Write an equation relating the forces at the top of the hill to the centripetal acceleration. Hence find the speed at which the car just loses contact with the road.
  \par\hfill\textit{[4 marks]}
  \vspace{4.5cm}
  \end{minipage}

  \item \begin{minipage}[t]{\linewidth}
  At a speed of $20\ \text{m s}^{-1}$, calculate the normal contact force on the car at the top of the hill.
  \par\hfill\textit{[2 marks]}
  \vspace{3.5cm}
  \end{minipage}
\end{enumerate}

\begin{minipage}{\linewidth}
\textbf{Q7.} A conical pendulum consists of a mass of $0.25\ \text{kg}$ on a string of length $0.60\ \text{m}$, rotating so that the string makes an angle of $30°$ with the vertical.
\end{minipage}

\begin{enumerate}[label=(\alph*)]
  \item \begin{minipage}[t]{\linewidth}
  Calculate the radius of the circular path.
  \par\hfill\textit{[1 mark]}
  \vspace{2.0cm}
  \end{minipage}

  \item \begin{minipage}[t]{\linewidth}
  Calculate the tension in the string.
  \par\hfill\textit{[2 marks]}
  \vspace{3.0cm}
  \end{minipage}

  \item \begin{minipage}[t]{\linewidth}
  Calculate the angular speed of the mass.
  \par\hfill\textit{[3 marks]}
  \vspace{4.0cm}
  \end{minipage}
\end{enumerate}

\end{tcolorbox}

% ===== SECTION 10: MARK SCHEME =====
\section{Mark Scheme and Answers}

\begin{markscheme}

\textbf{Q1.} The radian is the angle subtended at the centre of a circle by an arc whose length is equal to the radius of the circle [1].

\vspace{0.5em}\textbf{Q2.} (a) $\omega = 120 \times 2\pi/60 = \mathbf{4\pi = 12.6\ \text{rad s}^{-1}}$ [2]. (b) $v = r\omega = 0.35 \times 12.6 = \mathbf{4.4\ \text{m s}^{-1}}$ [1].

\vspace{0.5em}\textbf{Q3.} The direction of the velocity changes continuously [1]; acceleration is the rate of change of velocity — a change in direction (even at constant speed) constitutes acceleration, directed towards the centre of the circle [1].

\vspace{0.5em}\textbf{Q4.} $\omega = 3.0 \times 2\pi = 18.85\ \text{rad s}^{-1}$ [1]; $T = F = mr\omega^2 = 0.15 \times 0.80 \times 18.85^2$ [1] $= \mathbf{42.6\ \text{N}}$ [1].

\vspace{0.5em}\textbf{Q5.} The centripetal force is always perpendicular to the velocity (direction of motion) [1]; work done $= F\cos 90° \times d = 0$ — a perpendicular force does no work [1].

\vspace{0.5em}\textbf{Q6(a).} Diagram showing weight $mg$ downward and normal force $N$ upward, with $mg > N$ at speed [2].

\vspace{0.5em}\textbf{Q6(b).} At top of hill, net downward force provides centripetal force: $mg - N = mv^2/r$ [1]; at the point of losing contact, $N = 0$: $mg = mv^2/r$ [1]; $v = \sqrt{gr} = \sqrt{9.81 \times 120}$ [1] $= \mathbf{34.3\ \text{m s}^{-1}}$ [1].

\vspace{0.5em}\textbf{Q6(c).} $mg - N = mv^2/r$; $N = m(g - v^2/r) = 950(9.81 - 20^2/120) = 950(9.81 - 3.33)$ [1] $= 950 \times 6.48 = \mathbf{6160\ \text{N}}$ [1].

\vspace{0.5em}\textbf{Q7(a).} $r = L\sin\theta = 0.60\sin 30° = \mathbf{0.30\ \text{m}}$ [1].

\vspace{0.5em}\textbf{Q7(b).} $T\cos\theta = mg$; $T = mg/\cos 30° = (0.25 \times 9.81)/\cos 30°$ [1] $= \mathbf{2.83\ \text{N}}$ [1].

\vspace{0.5em}\textbf{Q7(c).} $T\sin\theta = mr\omega^2$ [1]; $\omega^2 = T\sin\theta/(mr) = 2.83\times\sin 30°/(0.25\times 0.30)$ [1]; $\omega = \sqrt{2.83\times0.5/0.075} = \sqrt{18.9} = \mathbf{4.34\ \text{rad s}^{-1}}$ [1].

\end{markscheme}

% ===== SECTION 11: REVISION CHECKLIST =====
\newpage
\section{Revision Checklist}

Use this checklist to track your understanding. Tick each box when you are confident:

\begin{tcolorbox}[colback=white, colframe=darkblue, breakable]
\renewcommand{\arraystretch}{1.5}
\begin{tabular}{@{} c p{10.5cm} r @{}}
\toprule
 & \textbf{Learning Objective} & \textbf{Confidence (1--3)} \\
\midrule
$\square$ & Define the radian and convert between degrees and radians & \\
$\square$ & Define angular speed $\omega$ and use \nf{\omega = 2\pi/T = 2\pi f} & \\
$\square$ & Use \nf{v = r\omega} to relate linear and angular speed & \\
$\square$ & Explain why an object in uniform circular motion is accelerating & \\
$\square$ & State that centripetal acceleration is directed towards the centre & \\
$\square$ & Use \nf{a = r\omega^2} and \nf{a = v^2/r} & \\
$\square$ & Use \nf{F = mr\omega^2} and \nf{F = mv^2/r} & \\
$\square$ & Identify the physical force providing centripetal force in a given situation & \\
$\square$ & Explain why centripetal force does no work & \\
$\square$ & Solve problems involving objects at the top/bottom of vertical circles & \\
$\square$ & Solve conical pendulum problems by resolving tension components & \\
\bottomrule
\end{tabular}
\vspace{0.5em}

\textit{Key: 1 = Need more work \quad 2 = Getting there \quad 3 = Confident}
\end{tcolorbox}

\vspace{1em}
\begin{motivationbox}
\centering
\textbf{\color{darkblue}Good luck with your revision!}\\[0.2em]
{\color{darkgray}\small Circular motion is the bridge between mechanics and fields — once you are confident with centripetal force and acceleration, gravitational and electric field problems fall into place naturally. Always ask: what real force is pointing towards the centre?}
\end{motivationbox}
